\section{Population eigenfunctions and the required core-growth class}
\label{pm:section}

An eigenfunction equation requires a growth class to characterize the
core exponent.  The following explicit family shows why this condition
cannot be omitted, even if the full eliminated kernel and its stopping
boundary are handled exactly.

Write $P_t^\infty$ and $\mathcal L_\infty$ for the semigroup and
proper-birth generator of the actual one-stub inert process.  If $T$ is
the first core-clock ring, with a null choice of the stub retained as
a ring, its discounted eliminated kernel is
\[
 (Q_sf)(y)=\mathbb E_y[e^{-sT}f(Y_T)].
\]
Core rings and leaf-clock rings retain their original source marks.

\begin{theorem}[Positive population modes at every discount]
\label{pm:all-discounts}
Consider the actual one-stub inert process.  Write $N(y)=n$ for its
real population and $D_0(y)$ for its root degree including the stub.
For every real $s>0$ there is a strictly positive function $h_s$ with
\[
 \mathcal L_\infty h_s=s h_s,
 \qquad P_t^\infty h_s=e^{st}h_s,
 \qquad Q_s h_s=h_s.
\]
Here $Q_s$ is the actual discounted first-core-ring kernel, including
the root's null stub choice.  The functions satisfy
\begin{equation}\label{pm:gamma-bounds}
 \frac{\Gamma(n-\tfrac12+s)}{\Gamma(n-\tfrac12)}
 \le h_s(y)\le
 \frac{\Gamma(n+s)}{\Gamma(n)}.
\end{equation}
In fact $h_s(z)>h_s(y)$ at every proper birth $y\to z$.  Under the
corresponding Doob transform, if $T$ is the first core ring and
$k=C(y)$, then
\[
 \Pr_y^{h_s}(T>t)\le e^{-kt}.
\]
Thus these modes have no core-ring escape at any positive discount.
\end{theorem}
\begin{proof}
After removing null rings, the proper total birth rate is
\[
 \lambda(y)=n-\frac1{D_0(y)},\qquad n-\tfrac12\le\lambda(y)\le n.
\]
The process $e^{-t}N(Y_t)$ is a nonnegative supermartingale: its
compensator drift is $-e^{-t}/D_0(Y_t)$.  It therefore has an almost-sure
limit $W_y$.

Use common independent unit exponentials for the holding times at each
population.  This sandwiches $N(Y_t)$ between pure-birth chains with
rates $n-1/2$ and $n$, started from the same population.  For a chain
with rate $a+j$ after $j$ jumps, the jump count at time $t$ has
negative-binomial generating function
\[
 \left[\frac{e^{-t}}{1-(1-e^{-t})z}\right]^a.
\]
Consequently its rescaled population converges to a gamma random
variable of shape $a$ and rate one.  This also follows from the
positive martingale $e^{-t}(a+j(t))$; its moments of every fixed
positive order are bounded.  In the above coupling this gives
\[
 0<W_y^-\le W_y\le W_y^+\quad\hbox{almost surely},
 \qquad W_y^-\sim\Gamma(n-\tfrac12,1),\quad
 W_y^+\sim\Gamma(n,1).
\]
No independence between these coupled variables is asserted.
The upper comparison also gives uniform integrability of every fixed
power of $e^{-t}N(Y_t)$, by taking a larger moment.

Define $h_s(y)=\mathbb E_y W_y^s$.  The gamma moment formula proves
\eqref{pm:gamma-bounds}.  The Markov property and the rescaling of the
future population limit give
$h_s(y)=e^{-st}\mathbb E_y h_s(Y_t)$, proving the semigroup equation.
Alternatively, conditioning on the first proper birth gives the finite
identity
\[
 h_s(y)=\frac1{s+\lambda(y)}\sum_z b(y,z)h_s(z),
\]
which is the generator equation.

The function $x\mapsto\Gamma(x+s)/\Gamma(x)$ is strictly increasing
for $s>0$: couple independent gamma variables of shapes $x$ and a
positive increment, and use strict increase of their $s$th powers.
A successor has population $n+1$, so its lower bound in
\eqref{pm:gamma-bounds} is strictly larger than the upper bound for
its predecessor.  Thus $h_s(z)>h_s(y)$ at every proper birth.

Keep the original source-clock marks in the Doob process.  Every
proper marked rate is multiplied by $h_s(z)/h_s(y)>1$, while a null
stub choice retains its rate.  The transformed proper total rate is
$\lambda(y)+s\le N(y)+s$, so the process is nonexplosive.  Its
core-ring hazard is at least its original value $k$.  Until the first
core ring, $k$ is constant, proving the displayed survival bound.
The finite-time density ratio is $e^{-st}h_s(Y_t)/h_s(y)$.
Stopping this martingale at $T\wedge t$ and letting $t\to\infty$
therefore gives $Q_s h_s=h_s$.  More precisely the omitted stopping
boundary is at most $h_s(y)e^{-kt}$.
\end{proof}

There is an especially transparent version in the original free-tree
generator $\mathcal L$, whose total birth rate is exactly $n$:
\[
 H_s(T)=\frac{(s+2)_{N(T)-2}}{(N(T)-1)!}
       =\prod_{j=2}^{N(T)-1}\left(1+\frac{s}{j}\right).
\]
It satisfies $\mathcal L H_s=sH_s$ for every $s>0$, since
$H_s(n+1)/H_s(n)=(n+s)/n$.  At rational $s$, every state value is
rational; at algebraic $s$, every state value is algebraic.  These
are population modes of the actual full model.  Neither positivity,
algebraic state values, nor an exact generator eigenfunction equation
alone forces $s$ to be the core-growth exponent.

The new modes do not resolve the stopping boundary of the previously
constructed core-reward-normalized critical eigenfunction.  They are
a different family, constructed from the population limit.  Moreover
they are excluded from uniform comparability with $C$: reachable
one-core stars have $C=1$ and arbitrarily large population, whereas
\eqref{pm:gamma-bounds} grows as $n^s$.

\subsection{An exact variational equation in the core-growth class}

For a finite inert state $y$, define
\[
 V_s(y)=\int_0^\infty e^{-st}(P_t^\infty C)(y)\,dt.
\]
For the uniform newborn resolvent, let $x$ range over the finite
admissible newborn-core environments of the original manuscript,
including their permitted inert-boundary completions.  Let $F_x(t)$
be the expected number of cores in the tagged newborn's descendant
family, including the newborn itself, and put
\[
 W(s)=\sup_x\int_0^\infty e^{-st}F_x(t)\,dt.
\]
We use the canonical-resolvent theorem of the original
manuscript\cite{Long}: for $s>\mu$, $W(s)<\infty$ and
\[
 C/s\le V_s\le(1+W(s))C/s,
 \qquad (s-\mathcal L_\infty)V_s=C.
\]
The supremum in $W(s)$ remains outside the integral.

Let $\mathcal A_C$ consist of functions on the reachable inert states
for which $C\le f\le K_f C$ with one finite constant $K_f$ valid
at every state.  All generator sums below are finite pointwise.

\begin{theorem}[Core-growth variational characterization]
\label{pm:variational}
The original exponent satisfies
\[
 \boxed{\displaystyle
 \mu=\inf_{f\in\mathcal A_C}\sup_y
                  \frac{(\mathcal L_\infty f)(y)}{f(y)}.}
\]
The infimum is not asserted to be attained.  If a positive eigenfunction
in this comparability class exists at a positive parameter $s$, then
$s=\mu$.
\end{theorem}
\begin{proof}
If $\mathcal L_\infty f\le\alpha f$ and $f\ge C$, localization
at a finite population and Fatou's lemma give
$P_t^\infty f\le e^{\alpha t}f$.  Nonexplosion removes the stopping
time.  Thus $F_\infty(t)\le e^{\alpha t}f(y_0)$, and the established
mean exponent gives $\mu\le\alpha$.

Conversely, for any $s>\mu$, the existing canonical resolvent theorem
gives
\[
 \frac{C}{s}\le V_s\le\frac{1+W(s)}s C,
 \qquad (s-\mathcal L_\infty)V_s=C.
\]
Thus $f=sV_s$ belongs to $\mathcal A_C$ and
$\mathcal L_\infty f=sf-sC\le sf$.  Taking $s\downarrow\mu$
proves the formula.  For an eigenfunction in this class, stopping at
population $M$ and using $f(Y_{t\wedge\tau_M})\le K_fN(Y_t)$
give uniform integrability at each fixed time.  Thus localization
justifies the semigroup equality $P_t^\infty f=e^{st}f$;
equivalently its Doob total rate
is $\lambda+s$, so the transform is conservative.  Comparability
with $C$ then identifies its growth exponent with $s=\mu$.
\end{proof}

\begin{corollary}[Uniform resolvent blow-up at the critical point]
\label{pm:uniform-pole}
For $s>\mu$, put
\[
 K(s)=\sup_y\frac{V_s(y)}{C(y)}<\infty.
\]
Then
\[
 \boxed{\displaystyle
 K(s)\ge\frac1{s-\mu},\qquad
 W(s)\ge\frac{\mu}{s-\mu}.}
\]
In particular $K(\mu)=W(\mu)=\infty$, in the extended-value sense,
and $W(s)<\infty$ holds exactly for $s>\mu$.
\end{corollary}
\begin{proof}
Since $V_s\le K(s)C$,
\[
 \mathcal L_\infty V_s=sV_s-C
 \le\left(s-\frac1{K(s)}\right)V_s.
\]
The same growth comparison gives $\mu\le s-1/K(s)$.
Combine this with $K(s)\le(1+W(s))/s$ to obtain both bounds.
Monotonicity in the discount gives divergence of the suprema at
$s=\mu$; below $\mu$ the initial inert resolvent already diverges.
\end{proof}

This is a pole-type lower bound for a \emph{uniform supremum} of
resolvents.  It does not show that the resolvent at any specified finite
initial state diverges at $\mu$, and it does not establish a pole of
the scalar free-tree transform $\mathcal M$.  No interchange of the
state supremum with a time integral is used.  Given a certified lower
bound $a\le\mu$ and $s>\mu$, it also gives the rational lower bounds
$K(s)\ge1/(s-a)$ and $W(s)\ge a/(s-a)$ when $s,a$ are rational.
